\begin{abstract}
Conventional Hyperspectral Imaging (HSI) is based on acquiring either reflectance, emission, or excitation profiles from each pixel of an image.
Emission-based HSI can be expanded into a fourth dimension by varying the excitation wavelength, which enables an improved discrimination of materials that contain mixed fluorophores.
Multi-excitation HSI (ME-HSI) creates a 4D dataset encompassing signal intensity across two spatial and two spectral dimensions, the latter being the emission and excitation wavelengths.
The processing of ME-HSI datasets presents notable challenges due to the numerous spectral layers per pixel, introducing considerable redundancy and computational burden.
Existing band selection methods, developed for 3D HSI, fail to capture the nonlinear cross-excitation correlations present in ME-HSI.
To fill this gap, we developed a deep learning framework comprising three stages: (1) a 3D convolutional autoencoder with parallel excitation branches that learns a unified latent representation of ME-HSI data, (2) a perturbation-based attribution method that traces reconstruction sensitivity to individual excitation-emission wavelength pairs, and (3) Maximum Marginal Relevance (MMR) selection that balances informativeness with spectral diversity.
The framework is evaluated on two ME-HSI fluorescence datasets with distinct sample chemistries: a four-class lichen autofluorescence dataset (3{,}072 configurations sweep) and a three-class collagen sponges dataset that tests transfer to a chemically distinct fluorophore family.
On lichens, 95.2\% accuracy is achieved with 80 of 192 bands (58\% data reduction, 7.0 points above the 88.2\% baseline), and 89.4\% is maintained at only 9 bands (95\% reduction).
On collagen sponges, the same pipeline---unchanged---reaches 85.6\% accuracy at 30 of 158 bands (81\% reduction, 5.8 points above baseline); a ten-classifier evaluation confirms the selected bands are intrinsically informative rather than tuned to a single classifier.
Robustness analysis on the Lichens dataset further validates the learned selection against $10{,}000$ random band combinations: the random distribution centers at $46.1\%$ accuracy (median $44.9\%$), while the learned 13-band selection scores $90.2\%$, well above the random tail.
The framework facilitates practical applications of ME-HSI by reducing acquisition complexity while improving discriminative power.
\end{abstract}
